\section{Hướng tăng tốc tổng quát bằng mảng systolic/GEMM}
ConvAU của Kiningham, Graczyk và Ramkumar đại diện cho hướng lấy nhân ma trận đặc làm trung tâm. Thiết kế dùng phép nhân 8 bit và tích lũy 32 bit, giữ trọng số tại PE, cấp đầu vào qua FIFO SRAM và lưu kết quả trung gian trong hệ thống đệm. Host nạp trọng số cùng metadata lượng tử hóa; phần cứng đảm nhiệm tính toán ma trận, lượng tử hóa và ReLU. Báo cáo không xác định một lõi RISC-V hoặc giao thức MMIO/RoCC tương ứng, nên không thể xem ConvAU là một SoC cùng cấu hình với luận văn~\cite[Sections~3--4]{kiningham2016design}.

\subsection{Thông lượng, chuẩn bị dữ liệu và mức sử dụng mảng}
Dòng MAC đều đặn cho phép tái sử dụng toán hạng qua các PE và khai thác song song khi kích thước ma trận phù hợp. Tuy nhiên, tích chập cần được ánh xạ sang phép nhân ma trận. ConvAU trình bày phép biến đổi im2col để sắp cửa sổ đầu vào thành cột~\cite[Section~4.2]{kiningham2016design}. Nếu vật chất hóa ma trận đó, các phần tử cửa sổ chồng lấn bị lặp, làm tăng lưu trữ và lưu lượng; nếu sinh dữ liệu theo luồng, chi phí chuyển sang bộ tạo địa chỉ và bộ đệm. Vì vậy, so sánh phải tính cả chuẩn bị dữ liệu, nạp trọng số, làm đầy và tháo cạn đường ống.

Đối với workload của luận văn, pointwise có cấu trúc trộn kênh phù hợp phép nhân ma trận, còn depthwise tính độc lập từng kênh~\cite[Section~3.1]{howard2017mobilenets}. Từ cấu trúc đó có thể suy ra rằng ánh xạ depthwise hoặc ma trận nhỏ lên mảng đặc lớn cần lịch ghép tác vụ hay chia khối phù hợp để tránh nhiều PE không có việc hữu ích. Đây là nhận xét về điều kiện sử dụng mảng, không phải số đo mức sử dụng của ConvAU trên TinyMobileNet và cũng không có nghĩa mọi mảng systolic đều xử lý depthwise kém.

\subsection{Phạm vi đánh giá và bài học cho hệ thống hiện tại}
ConvAU được đánh giá bằng RTL, tổng hợp và luồng ASIC, với SqueezeNet lượng tử hóa làm workload~\cite[Section~5]{kiningham2016design}. Thông lượng và công suất trong báo cáo thuộc điều kiện đó; không dùng chúng để suy ra ưu thế năng lượng hoặc độ trễ của FPGA trong luận văn. So sánh số bit cũng chỉ là một tiêu chí: cùng độ rộng tensor và tích lũy chưa bảo đảm quy tắc làm tròn hay saturation giống nhau.

Khác với ConvAU, MyTinyMobileNetSoC không thực hiện mạng bằng một mảng ma trận đặc chung và không dựng ma trận im2col đầy đủ trong đường chạy hiện tại. Dữ liệu được chuẩn bị theo patch hoặc vector kênh cho ngoại vi chuyên biệt. Lựa chọn này phù hợp phạm vi hiện thực TinyMobileNet đã xác định, nhưng đánh đổi bằng lịch firmware phụ thuộc shape và phần tính toán còn ở CPU. Bài học từ hướng systolic là phải cân đối số MAC với khả năng nuôi dữ liệu và miền tái sử dụng, thay vì xem kích thước mảng là thước đo đủ cho hiệu năng hệ thống.
